\section{Ordinary graph branes and a simultaneous algebraic mirror}
\label{mirror:section}

We now work at the ordinary point $t=0$, retaining the integer grading
on the graph category $\mathcal D_\Lambda$.  The purpose of this section
is to compare this directed category with actual line bundles.  The
comparison is first made on finite lists over the complete Novikov
field.  Finite algebraic models then assemble those comparisons after
one further field extension.  Only the ordinary graph branes enter
family Floer theory; the module constructed from $L$ will enter through
the algebraic localization in the next section.

Put
\[
 Y_\Lambda=E(\qnov)^8,\qquad
 H_Y=\bigotimes_{a=1}^8 H_{Y,a},
\]
where $E(\qnov)$ is the Tate elliptic curve with multiplicative period
$\qnov=\qnov^1$, and $H_{Y,a}$ is a principal theta line from the $a$th
factor.  Write
\[
 \Sigma=\{h_0k\,1_8:k\geq1\}\cup\{S_-\},\qquad
 N_\sigma=\bigotimes_{a=1}^8H_{Y,a}^{-\sigma_{aa}}.
\]
Thus $N_{h_0k1_8}=H_Y^{-h_0k}$ and
$N_{S_-}=\bigotimes_{a\leq4}H_{Y,a}^{-1}
\otimes\bigotimes_{a>4}H_{Y,a}$.  A \emph{slope type} means an element
of $\Sigma$; the many labels of one slope type are its Hamiltonian
perturbations in the successive layers.

\subsection{The finite ordinary comparison}

We use the following precise scope of established family Floer theory.
For a finite transverse collection of embedded, integer-graded,
tautologically unobstructed branes in a closed nonsingular Lagrangian
torus fibration with $\pi_2$ of the base zero, compatible relative
Pin data give a faithful $\Ainf$ functor to twisted analytic perfect
complexes, after choosing a sufficiently fine family cover
\cite[Theorem~2.10, \S2.4.1, and Appendix~A]{abouzaid-faithful}.
The corresponding cohomology maps are also surjective
\cite[Theorem~1.1 and Corollary~1.2]{abouzaid-hms}.  The finite
$\Ainf$ construction, rather than only its map on morphism cohomology,
is the input used here.

All its hypotheses hold for our graphs.  The fibration
\[
 T\longrightarrow \R^8_q/\Z^8
\]
has closed total space and no singular fibers.  Both its base and total
space are aspherical.  For
$X=\{p=A_Xq+\nabla f_X(q)\}$, the map on fundamental groups is
$b\mapsto(b,A_Xb)$, hence is injective.  The relative homotopy exact
sequence gives $\pi_2(T,X)=0$.  Consequently every disk on $X$ has
area zero in its relative class, while every nonconstant tame
holomorphic disk has positive area.  The same argument excludes
spheres.  These branes are therefore tautologically unobstructed for
every tame almost complex structure, with zero bounding cochain.
Their spin structures come from the base frame, their local systems
are trivial and valuation-unitary, and their integer graph gradings
are obtained by continuation from constant symmetric graph planes.
The fixed squared complex volume form represents the fibration
grading: its phase on the constant vertical frame is constant.

The two possible twists are trivial with these choices.  The global
zero section may be used in every family chart, so all differences of
local sections and their primitive cocycle are zero.  The global
base frame trivializes the ordinary-spin comparison cocycle as well;
the relative background is zero.  These are the section and base-spin
terms distinguished in \cite[Remark~2.16]{abouzaid-hms}.  The rigid dual
has uniformizing coordinates $z=(z_1,\ldots,z_8)$ and periods
$z\mapsto\qnov^b z$, $b\in\Z^8$, so it is $Y_\Lambda^{\mathrm{an}}$.
The preceding verification applies also to $C$: no extra $B$-weight is
being imposed in this section.

Choose the covariant convention for the mirror functor.  If a family
model uses the contravariant convention of
\cite[Equation~(1.1)]{abouzaid-hms}, compose once with perfect duality.
This changes all numerical signs together.  Restricting an ordinary
finite functor to strictly ordered words gives a nonunital functor on
the directed sector.  Every word in its boundary identities remains
strictly ordered.  Proposition~\ref{cat:ordinary-comparison} lets
us use the prescribed operations of $\mathcal D_\Lambda$: compare the
two finite sets of ordinary data on separately labeled copies, take
quasi-inverses, and then restrict to strict words.  This comparison
involves only ordinary weakly exact Floer theory.

\begin{lemma}[Numerical type of a graph image]
\label{mirror:graph-image}
For any such finite construction, the image of a graph $X$ is a line
bundle, up to an integer shift.  Before fixing the uniform convention
sign, its numerical type is $\varepsilon A_X$, with one
$\varepsilon\in\{1,-1\}$ for the whole construction.  More precisely,
it is a shift of
\[
 \bigotimes_{a=1}^8H_{Y,a}^{\varepsilon(A_X)_{aa}}\otimes P_X,
 \qquad P_X\in\Pic^0(Y_\Lambda).
\]
\end{lemma}

\begin{proof}
A graph meets every torus fiber once and transversely.  Its local
family complex is consequently free of rank one in one integer
degree, with zero differential.  We determine its transition
characters; this is the additional information needed to identify
the numerical line type.

On the real universal cover of the base, use
\[
 \ell(q)=A_Xq+\nabla f_X(q)=dg(q),\qquad
 g(q)=\tfrac12q^tA_Xq+f_X(q),
\]
and the fiber paths $s\ell(q)$ from the zero section.  In family
continuation, a coefficient carries the character of the fiber
boundary closed by these paths
\cite[\S\S2.2 and 8.2]{abouzaid-faithful}.  For compatible real lifts
this character is zero.  Indeed, lift the projection of a comparison
strip along its prescribed local base path.  A further lift of its
graph boundary lies in
$p=\ell(q)+a$ for one constant $a\in\Z^8$.  The endpoint paths cancel
this same offset at both ends, and the other boundary lies over the
prescribed base path.  Closing the strip therefore leaves no
additional fiber winding.  Equivalently, such a winding would be a
nonzero fiber loop bounding in
$\R^8\times(\R^8/\Z^8)$, contrary to injectivity of its fiber
fundamental group.

The family trivialization factor contributes no missing character.
Its exponent is $df_{ij}-dg_j+dg_i$ in
\cite[Equation~(8.16)]{abouzaid-faithful}.  Our common zero section
gives $f_{ij}=0$, and the common real lift gives $dg_i=dg_j$.
One can make the calculation for the undisplaced section, already
transverse to every fiber, or use continuation of small perturbations.
For the latter, local identifications from $q$ to $q+v$ change the
angle by $\ell(q+v)-\ell(q)$ and preserve this lifted-path calculation.
In absolute multiplicative coordinates, every term of a local
transition thus has character exponent zero.  Its sum is a scalar
in $\Lambda$, nonzero by continuation of the one-generator
cohomology.

All local generator degrees agree on the real cover.  A homotopy
between the degree-zero transition maps has negative degree and is
zero; the higher descent maps vanish for the same reason.  The
scalar transitions form an ordinary locally constant cocycle.
Contractibility of the real cover permits scalar rescaling of the
generators to trivialize this cocycle.  The section and frame also
trivialize the twisting data there.  Hence no invertible analytic
function is left undetermined in the following period calculation.

A base period $b\in\Z^8$ changes the angle lift by $A_Xb$.
In the frame just chosen, the period multiplier is therefore
\begin{equation}
 j_X(b,z)=c_bz^{-\varepsilon A_Xb},\qquad c_b\in\Lambda^*.
 \label{mirror:period-character}
\end{equation}
The scalar is independent of the local polytope because the
transition frames agree and the real base is connected.  The
multiplier cocycle identity says
\[
 c_{b+c}=c_bc_c\qnov^{-\varepsilon b^tA_Xc}.
\]
For a diagonal integral symmetric matrix $A$, the product theta line
of type $\varepsilon A$ has multiplier
\[
 j_{\varepsilon A}(b,z)
   =\qnov^{-\varepsilon b^tAb/2}z^{-\varepsilon Ab}.
\]
Dividing Equation~\eqref{mirror:period-character} by this multiplier
leaves the scalar character
$d_b=c_b\qnov^{\varepsilon b^tA_Xb/2}$, for which
$d_{b+c}=d_bd_c$.  Such a character defines a line in $\Pic^0$:
writing $d_b=\prod_a d_{e_a}^{b_a}$, it is the quotient of the
principal product theta line translated by
$(d_{e_a}^{-1})_a$ and the original theta line.  Thus the analytic
image has exactly the asserted type.  Projective rigid GAGA
algebraizes this line over $\Lambda$; its applicable form over a
complete, possibly nondiscretely valued field is recalled in
\cite[p.~502]{conrad-gaga}.
\end{proof}

\begin{lemma}[Normalization of signs, shifts, and copies]
\label{mirror:normalization}
Suppose a finite list contains three distinct layer-zero scalar
slopes and the layer-zero representative of every slope type on the
list.  After a common shift its images are unshifted lines
\begin{equation}
 L_X=N_{A_X}\otimes P_{A_X},\qquad
 P_{A_X}\in\Pic^0(Y_\Lambda),
 \label{mirror:normalized-lines}
\end{equation}
with one degree-zero factor per slope type.  The finite directed
functor is an isomorphism on cohomology for every strict arrow.
\end{lemma}

\begin{proof}
For distinct slope types the matrix difference is definite.  Using
the small layer-zero perturbations, each lifted sheet pair intersects
once; there are $|\det(A_X-A_{X'})|$ projected intersections, all of
one degree.  The differential is therefore zero.  Hamiltonian
continuation gives the same cohomology for any copies.  The total
Ext dimension of the corresponding lines is the same determinant:
their quotient is an external product of nonzero-degree elliptic
lines, with an arbitrary $\Pic^0$ factor, and elliptic
Riemann--Roch and K\"unneth give exactly this dimension.  Faithfulness
is consequently fullness for these pairs.

Two copies of the same slope have the self-Floer cohomology of an
eight-torus.  Their images before shifts have degree-zero quotient.
A nontrivial such quotient has zero cohomology: on the product of
elliptic curves a $\Pic^0$ line is an external product, and at least
one nontrivial degree-zero factor has both cohomology groups zero.
Faithfulness therefore forces the quotient to be trivial.  Its Ext
groups have the torus Betti dimensions, so faithfulness is again
fullness.  We may identify all copies with the same line, leaving
their shifts to determine.

Write a layer-zero image as $L_i[s_i]$.  For $i>j$ the source Hom
from $G_i$ to $G_j$ is nonzero only in degree zero.  If
$\varepsilon=1$ in Lemma~\ref{mirror:graph-image}, its line quotient
is anti-ample and has Ext only in degree eight.  Thus
$s_j-s_i=8$ for every such pair.  Three slopes contradict the
additivity of shift differences.  Hence $\varepsilon=-1$.
The quotients are now ample and have Ext only in degree zero, so
all scalar shifts agree.  The layer-zero arrows $G_j\to C$ also
have degree zero.  Their line quotients have the positive diagonal
type $h_0j\,1_8-S_-$, so the shift for $C$ is the same.
Finally, the source cohomology between copies of a slope has support
in $[0,8]$, with both endpoints nonzero.  Since their underlying
lines agree, fullness forces their shift difference to vanish.
A common shift of the entire construction now gives
Equation~\eqref{mirror:normalized-lines}.
\end{proof}

\subsection{Algebraic enhancements and uniform generation}

We will need maps into chosen algebraic models, including their
higher components.  Here is a concrete enhancement of the GAGA
comparison just used.  Take a finite affine algebraic cover
$\mathfrak U$ of $Y_\Lambda$, an affinoid family cover on the analytic
side, and a common finite admissible affinoid refinement
$\mathfrak V$ subordinate to both.  Such a refinement exists by
quasi-compactness.  All algebraic and analytic intersections used
here are affine and affinoid, respectively, by separatedness.

For honest bundles $P,Q$, first form \v Cech Hom complexes using all
index tuples, with Alexander--Whitney cup composition.  Assign each
refinement open one containing open from each original cover.
Analytification followed by restriction gives actual cochain maps
\[
 \check C(\mathfrak U,P^\vee\otimes Q)
 \longrightarrow
 \check C(\mathfrak V,(P^\vee\otimes Q)^{\mathrm{an}})
 \longleftarrow
 \check C(\mathfrak U_{\mathrm{an}},
                         (P^\vee\otimes Q)^{\mathrm{an}}).
\]
The all-tuple convention allows arbitrary assignments of containing
opens.  These maps commute with differentials and cup composition.
Acyclicity of affine and affinoid intersections and coherent
cohomology GAGA make them quasi-isomorphisms.  Projection to strictly
increasing tuple complexes is also a cup dg map and a
quasi-isomorphism, giving the usual bounded \v Cech models.  Thus
this is an enhanced comparison on lines and bounded bundle
complexes.  Passing to finite cones and retracts gives the comparison
on their split-closed hulls.  In the family-cover category, the
rank-one images of Lemma~\ref{mirror:graph-image} have ordinary strict
descent, so analytic descent identifies their Homs and compositions
with exactly this analytic \v Cech model.  No algebraization of an
arbitrary analytic gerbe is involved.

\begin{lemma}[Nine lines generate uniformly]
\label{mirror:generation}
There are integers $1\leq k_1<\cdots<k_9$ such that, for every
choice $P_j\in\Pic^0(Y_\Lambda)$, the nine lines
\[
 H_Y^{-h_0k_j}\otimes P_j,\qquad 1\leq j\leq9,
\]
split-generate $\Perf(Y_\Lambda)$.  The same assertion holds after
any algebraically closed field extension of $\Lambda$.
\end{lemma}

\begin{proof}
Choose a vector bundle $V$ which classically generates
$\Perf(Y_\Lambda)$.  The classical-generator theorem
\cite{bondal-vdb} gives a perfect generator; a bounded vector bundle
resolution and the direct sum of its terms give such a $V$.
We construct nine evaluation steps uniformly in the degree-zero
factors.

Put $K_0=V$ and $S_j=(\Pic^0Y_\Lambda)^j$.  Suppose $K_{j-1}$ is
a vector bundle on $Y_\Lambda\times S_{j-1}$.  Pull it to
$Y_\Lambda\times S_j$, tensor by the inverse Poincar\'e line from
the last factor and by $H_Y^{h_0k_j}$.  Relative Serre vanishing and
generation for the projection to the proper parameter scheme $S_j$
give, for one sufficiently large $k_j>k_{j-1}$, a locally free
pushforward $W_j$, formation of which commutes with base change,
and an evaluation surjection.  Tensoring back gives
\[
 0\longrightarrow K_j\longrightarrow
 W_j\otimes(H_Y^{-h_0k_j}\otimes\mathcal P_j)
 \longrightarrow K_{j-1}\longrightarrow0.
\]
Here $W_j$ is pulled back from $S_j$, and the kernel is locally free
because the quotient is locally free.  This completes the induction
with exponents independent of the points $P_j$.

At any nine parameter points the resulting exact sequence is
\[
 0\to K_9\to W_9\otimes L_9\to\cdots\to
 W_1\otimes L_1\to V\to0,
 \qquad L_j=H_Y^{-h_0k_j}\otimes P_j,
\]
where now the $W_j$ are finite-dimensional vector spaces.  If $C_V$
is the complex of the nine displayed line-bundle terms, there is a
triangle
\[
 K_9[8]\longrightarrow C_V\longrightarrow V
       \longrightarrow K_9[9].
\]
Since $V$ is locally free and $\dim Y_\Lambda=8$,
\[
 \Ext^9(V,K_9)=H^9(Y_\Lambda,V^\vee\otimes K_9)=0.
\]
The triangle splits, so $V$ is a summand of $C_V$.  The nine lines
therefore split-generate.  All the evaluation sequences and their
base-change properties were constructed over the entire parameter
schemes.  They and the splitting argument remain valid after the
stated field extensions.  The generator also remains a generator:
by adjunction, a quasi-coherent complex right-orthogonal to its
pullback has zero pushforward to $Y_\Lambda$; the affine projection
is conservative on quasi-coherent complexes, so that complex is zero.
\end{proof}

Fix these nine layer-zero labels as a core.  Enlarge any finite list
to contain the core and the layer-zero representatives of all its
slope types.  Lemma~\ref{mirror:normalization} and the enhanced GAGA
comparison give a nonunital $\Ainf$ functor from its ordered sector
to an algebraic enhancement on the lines
$N_\sigma\otimes P_\sigma$, with isomorphisms on strict-arrow
cohomology.  The images split-generate $\Perf(Y_\Lambda)$ by
Lemma~\ref{mirror:generation}.  We may replace this enhancement by
any chosen quasi-equivalent line model, using minimal-model transfer
and quasi-inverses over the field \cite{seidel}.  Isomorphic copies
of a line may similarly be replaced by a single chosen line with
several labels.  This statement concerns the target enhancement:
it does not regard the finite directed source as a full ordinary
subcategory.

\subsection{Finite algebraic models over the Picard variety}

The preceding finite constructions need not agree when the list
grows.  We next arrange algebraic formulas for their possible targets
so that a simultaneous functor can be obtained without imposing
analytic compatibility.

Let $\widehat Y_\Lambda=\Pic^0(Y_\Lambda)$, with a fixed Poincar\'e
line.  Choose a finite affine open cover of $\widehat Y_\Lambda$.
For each slope type $\sigma$, the possible line is
$N_\sigma\otimes\mathcal P_\sigma$.  In a word containing several
labels of the same slope type, use \emph{one} parameter for that type.
Thus the parameter base for two distinct types is a product of two
chosen affine charts, whereas for two copies of the same type it is
just the one chart.

\begin{lemma}[Pairwise contractions with algebraic parameters]
\label{mirror:parameter-model}
For every pair of types and every choice of their affine parameter
charts, there is a contraction of its bounded algebraic \v Cech Hom
complex onto a finite projective graded cohomology module.  The
contractions can preserve the actual constant units.  Choosing them
pairwise and pulling them back to finite products of parameter
charts gives strictly unital $\Ainf$ line models whose operations at
each arity are algebraic maps of finite projective modules and
commute with arbitrary specialization.
\end{lemma}

\begin{proof}
Fix an ordered finite affine cover of $Y_\Lambda$.  Over an affine
parameter base $S=\Spec A$, let $C_{\sigma\tau}$ be the associated
strictly increasing \v Cech complex of
\[
 (N_\sigma\otimes\mathcal P_\sigma)^\vee
          \otimes(N_\tau\otimes\mathcal P_\tau)
\]
on $Y_\Lambda\times S$.  These pair complexes carry the usual cup
compositions after pullback to a common finite parameter product.
Their terms are projective $A$-modules, although generally of
infinite rank.  Indeed, on an affine intersection $U=\Spec B$ of
the chosen cover, sections of the line are a finite projective
module over $B\otimes_\Lambda A$.  This ring is free as an
$A$-module because $B$ is a $\Lambda$-vector space.  Such a section
module is therefore an $A$-direct summand of a free module.  Finite
sums preserve this property.  Affineness also shows that the terms,
differential, and cup composition commute with base change.

The cohomology modules are finite projective and commute with
specialization.  If $\sigma\ne\tau$, the quotient line has either
ample or anti-ample diagonal type: scalar types are totally ordered,
and $h_0k1_8-S_-$ is positive definite.  Its fiber cohomology is
concentrated in degree zero or eight, with constant dimension
$|\det(\sigma-\tau)|$, by elliptic Riemann--Roch and K\"unneth.
Proper coherent cohomology and base change then give the assertion.
For a repeated type the quotient is canonically $\cO_Y$, since its
two labels share the same parameter.  Its cohomology is simply
$H^*(Y_\Lambda,\cO_Y)\otimes_\Lambda A$.

We recall why a bounded complex of projective modules with
projective cohomology admits the required contraction.  Starting
in its highest degree, split the surjection from the top term to
its cohomology.  Its boundary submodule is then projective.  Split
the preceding differential onto that submodule; its kernel is
projective.  Split this kernel onto its projective cohomology, and
continue downward.  This decomposes the complex into its cohomology
and two-term identity complexes.  Hence there are $A$-linear maps
\[
 i:H(C)\to C,\qquad p:C\to H(C),\qquad h:C\to C[-1]
\]
with
\begin{equation}
 pi=\id,\qquad ip=\id-dh-hd,\qquad
 h^2=hi=ph=0.
 \label{mirror:contraction}
\end{equation}
For a repeated type use the same structure-sheaf contraction for
all its label pairs.  In degree zero choose the constant function
as the representative of $1\in H^0(Y,\cO_Y)$; the decomposition
can retain this summand.  Thus $h(1)=0$ and the units are preserved.

Choose these contractions for each pair and chart choice, once and
for all.  For a finite word pull them back from the relevant pair
parameter bases to the product indexed by its distinct types.
The identities in Equation~\eqref{mirror:contraction} survive every
pullback, so no flatness qualification is needed at specialization.
Homotopy transfer expresses the operation of arity $d$ as the finite
sum over rooted binary trees with $d$ leaves: insert $i$ at leaves,
cup composition at vertices, $h$ at internal edges, and $p$ at the
root, with the usual graded signs.  These are multilinear maps
between finite projective cohomology modules.  They are algebraic,
and every tree formula commutes with specialization.  The side
conditions and the preserved units give a strictly unital transfer.
Using a fixed pair contraction, independent of other types in a
word, makes these models compatible on all finite lists and on
lists with repeated labels.
\end{proof}

\subsection{A simultaneous ordered functor}

\begin{proposition}
\label{mirror:simultaneous}
There is an algebraically closed extension $\kappa_1/\kappa$ of
cardinality at most the continuum, a line
$L_\sigma=N_\sigma\otimes P_\sigma$ on
$Y_{\kappa_1}=Y_\Lambda\otimes_\Lambda\kappa_1$ for each slope
type, and a nonunital integer-graded $\Ainf$ functor
\begin{equation}
 J:\mathcal D_\Lambda\otimes_\Lambda\kappa_1
             \longrightarrow\mathcal B.
 \label{mirror:simultaneous-functor}
\end{equation}
Here $\mathcal B$ is a strictly unital algebraic $\Ainf$ model of
the full Hom category of these lines, with a separate label for
each graph copy.  The functor is a quasi-isomorphism on every strict
arrow.  It sends each $w_X:X\to X^+$ to a cohomological
isomorphism.  The split-closed hull of $\mathcal B$ is
$\Perf(Y_{\kappa_1})$.
\end{proposition}

\begin{proof}
Take an increasing exhaustion by finite lists containing the core
and the layer-zero representatives of their types.  Choose one
finite mirror functor on each list.  By the preceding normalization
it has one parameter point $p_{\sigma,n}\in\widehat Y_\Lambda$
per type; assign arbitrary points to absent types.  Choose one
affine chart containing each such point.

Fix a nonprincipal ultrafilter $\mathcal U$ on the positive integers
and set
\[
 \kappa_1=\kappa^{\N}/\mathcal U,
\]
the ordinary field ultrapower, with the constant embedding of
$\kappa$.  It is algebraically closed: roots of a fixed-degree
polynomial can be chosen in each coordinate wherever its leading
coefficient is nonzero.  Its cardinality is at most
$(2^{\aleph_0})^{\aleph_0}=2^{\aleph_0}$.  For each fixed type,
one of the finitely many chart choices occurs on a
$\mathcal U$-large set.  Its parameter sequence defines a
$\kappa_1$-point of that affine chart.  This uses only finitely many
coordinate functions and their algebraic relations.  Denote its
Poincar\'e line by $P_\sigma$.

At every finite stage transfer the target of the chosen functor to
the specialized model of Lemma~\ref{mirror:parameter-model}, using
the selected charts at that stage.  The enhanced comparison and
quasi-inverse transfer give actual higher functor components with
this target.  For a fixed word, its labels belong to all sufficiently
large lists, and all its chart choices agree with the selected
limiting charts on a $\mathcal U$-large set: only a finite
intersection of large sets is involved.  On that set, every target
operation is the specialization of one fixed algebraic module map.

Source Hom complexes for the word have fixed finite bases over
$\Lambda$.  Present each target finite projective as the image of
a finite idempotent matrix over its parameter ring.  Functor
components can then be written in finite matrices into those
images.  Taking their coefficient sequences gives maps over
$\kappa_1$ into the evaluated idempotent images.  For each arity
the $\Ainf$ functor equation has finitely many terms, and the target
tree formulas commute with this evaluation.  Its coordinatewise
validity on a large set therefore proves the equation over
$\kappa_1$.  Entries outside the relevant large set can be assigned
arbitrarily.

The first components remain quasi-isomorphisms.  To see explicitly
the finiteness of this assertion, fix a contraction of each source
complex to its finite-dimensional cohomology over $\Lambda$.
At every applicable stage, the first functor component induces an
invertible map from this fixed cohomology to the finite projective
target fiber.  Choose an inverse.  The two inverse equations are
finite matrix equations, including the idempotent equations if
the target is represented as a summand.  They pass to the
ultrapower, giving the claimed isomorphism on cohomology.

The resulting target is genuinely a full line model.  Base-changing
the pair contractions and performing their same finite tree
transfers computes all derived Homs and compositions of the chosen
lines on $Y_{\kappa_1}$, including arrows in both directions and
endomorphisms.  By Proposition~\ref{cat:continuation}, each $w_X$
was a nonzero degree-zero cohomology class.
Its image lies in the one-dimensional degree-zero Hom between two
copies of the same line, and remains nonzero by the strict-arrow
quasi-isomorphism.  It is consequently invertible.  Finally, the
nine core lines split-generate by Lemma~\ref{mirror:generation},
proving the hull assertion.
\end{proof}

\subsection{Extending the formal coefficients}

We make the extension in Proposition~\ref{mirror:simultaneous}
throughout the preceding algebraic constructions.  On the ordinary
fiber it is ordinary scalar extension by $\kappa\to\kappa_1$.
For Taylor coefficients it is the coefficientwise map
\begin{equation}
 R=\kappa[[t]]\longrightarrow R_1=\kappa_1[[t]].
 \label{mirror:complete-base-change}
\end{equation}
This is a local faithfully flat map.  Indeed $R$ is a DVR, $R_1$
is torsion-free over it and hence flat, and its residue extension
is the nonzero field extension $\kappa\to\kappa_1$.  In particular,
for every finite free periodic value complex $C$ constructed above,
\[
 H^r(C\otimes_RR_1)=H^r(C)\otimes_RR_1,
 \qquad
 (C\otimes_RR_1)/t=C/tC\otimes_\kappa\kappa_1.
\]
The same equalities hold for the chain actions and their induced
cohomology actions.  This uses the completed coefficient ring in
Equation~\eqref{mirror:complete-base-change}; no identification with
$\kappa_1\otimes_\kappa\kappa[[t]]$ is asserted.

Each graded theta-ring piece is finite free, so its base extension
has the same basis and multiplication formula with extended
coefficients.  Relative Proj and its grading line base-change, and
smoothness and projectivity persist.  The field-realization argument
continues to apply: the theta structure constants are unchanged,
and both $\kappa_1$ and $\operatorname{Frac}(R_1)$ have cardinality
at most the continuum.  All passage to the ultrapower and all
subsequent extension here concern finite algebraic matrices and
models.  We have not extended an analytic family Floer construction
to this field or evaluated a Floer series for $L$ there.

Henceforth rename $\kappa_1$ and $R_1$ as $\kappa$ and $R$, and
write $Y_\kappa=Y_\Lambda\otimes_\Lambda\kappa$.  With these
coefficients, Proposition~\ref{mirror:simultaneous} supplies the
simultaneous ordered map needed for localization.
